\documentclass[11pt,a4paper]{ltjsarticle}
\usepackage[margin=18mm]{geometry}
\usepackage{luatexja-fontspec}
\usepackage{xcolor}
\usepackage{hyperref}
\usepackage{enumitem}
\usepackage{tabularx}
\usepackage{tcolorbox}
\usepackage{tikz}
\usepackage{fancyhdr}
\usepackage{fontawesome5}
\usepackage{array}

\definecolor{ink}{HTML}{2B2119}
\definecolor{paper}{HTML}{FFF8E8}
\definecolor{wine}{HTML}{762F28}
\definecolor{gold}{HTML}{A8782C}
\definecolor{forest}{HTML}{315C42}
\definecolor{navy}{HTML}{31546B}
\definecolor{purple}{HTML}{653568}
\definecolor{muted}{HTML}{6B5B49}
\definecolor{line}{HTML}{C6A978}
\hypersetup{colorlinks=true,urlcolor=navy,linkcolor=wine}
\pagecolor{paper}
\color{ink}
\setlength{\parindent}{0pt}
\setlength{\parskip}{5pt}
\setlist{nosep,leftmargin=2em}
\pagestyle{fancy}
\fancyhf{}
\lhead{\small コード事件捜査局}
\rhead{\small PLAY GUIDE}
\cfoot{\small\thepage}
\renewcommand{\headrulewidth}{0.4pt}
\setmainjfont{Hiragino Sans}
\setsansjfont{Hiragino Sans}
\setmonojfont{Menlo}
\newfontfamily\titlejp{Hiragino Mincho ProN}

\newtcolorbox{casebox}[2][]{
  colback=white,colframe=#2,boxrule=0.8pt,arc=1mm,
  left=3mm,right=3mm,top=2.5mm,bottom=2.5mm,
  title={#1},fonttitle=\bfseries,coltitle=white,colbacktitle=#2
}
\newcommand{\taglabel}[2]{\tikz[baseline=(x.base)]\node[rounded corners=1pt,fill=#1,
text=white,font=\bfseries\small,inner xsep=6pt,inner ysep=2.5pt](x){#2};}
\newcommand{\key}[1]{\fcolorbox{line}{white}{\texttt{\strut #1}}}
\newcommand{\stepcircle}[1]{\tikz[baseline=(x.base)]\node[circle,fill=wine,text=white,
font=\bfseries,inner sep=2.2pt](x){#1};}

\begin{document}

\begin{titlepage}
\pagecolor{ink}
\color{white}
\vspace*{18mm}
{\color{gold}\rule{\linewidth}{1.2pt}}\par
\vspace{9mm}
{\titlejp\fontsize{28}{38}\selectfont コード事件捜査局\par}
\vspace{3mm}
{\fontsize{15}{22}\selectfont 機能紹介・遊び方・おすすめ攻略順\par}
\vspace{9mm}
{\color{gold}\rule{\linewidth}{1.2pt}}\par
\vfill
\begin{tcolorbox}[colback=white!7!ink,colframe=gold,coltext=white,arc=1mm,
boxrule=.8pt,left=5mm,right=5mm,top=5mm,bottom=5mm]
\large
プログラミングの勉強を「事件捜査」に変える学習ゲームです。\\
証拠を読み、問題を解き、未解決事件を片づけ、仲間と進捗を競いましょう。
\end{tcolorbox}
\vfill
{\large 最短ルートはこれ！\par}
\vspace{4mm}
\begin{center}
\taglabel{forest}{証拠を読む}
\quad$\rightarrow$\quad
\taglabel{wine}{推理に挑む}
\quad$\rightarrow$\quad
\taglabel{navy}{未解決を復習}
\quad$\rightarrow$\quad
\taglabel{purple}{ランキング}
\end{center}
\vfill
\small
Web版：\href{https://code-case-bureau-2026.web.app/}{https://code-case-bureau-2026.web.app/}\\
このガイドは2026年7月21日時点の公開版に対応しています。
\end{titlepage}
\pagecolor{paper}
\color{ink}

\section*{まず知っておきたいこと}

\begin{casebox}[このゲームの目的]{wine}
たくさんある既存の事件（学習章・問題）を進めながら、理解度、正答数、連続正解、
学習完了数などを伸ばします。ランキング専用の問題を毎日解く方式ではありません。
\textbf{普段の勉強が、そのままランキングの成績になります。}
\end{casebox}

\begin{tabularx}{\linewidth}{>{\bfseries}p{28mm}X}
事件一覧 & 章ごとの学習・問題へ入るメイン画面。進行度に応じて次の事件が開きます。\\
今日の一歩 & その日に取り組みやすい小さな目標。右上の「－／＋」で最小化できます。\\
捜査記録 & 総合評価、連続正解、正答率を確認できます。\\
自動保存 & 学習記録は端末に自動保存。ランキング部屋の記録はFirebaseで同期します。\\
\end{tabularx}

\section*{ホーム画面の見方}

\begin{tabularx}{\linewidth}{>{\bfseries}p{32mm}X}
事件一覧 & CASEを章ごとに選ぶメイン画面です。中央の事件カードから学習または演習を開始します。\\
ランキング & 友達とのランキング部屋、招待コード、スマホ・PCの進捗同期を管理します。\\
総合捜査 & 現在の条件で1000本ノックを開始・再開します。\\
未解決 & 間違えた問題や自分で登録した問題だけを解き直します。件数も表示されます。\\
捜査条件 & 1000本ノックの単元、重要度、難易度、問題数、快速モードを設定します。\\
証拠を読む／推理に挑む & 事件カードを押したとき、学習から入るか問題から入るかを切り替えます。\\
別解 & 表記が違っても意味が同じ答えを採用するか切り替えます。\\
観測を開く & 進捗、階級、実績バッジ、共有画像などを表示します。\\
\end{tabularx}

\begin{casebox}[スマホでの表示]{navy}
画面幅に合わせてサイドメニューやカードが縦並びになります。機能が消えるのではなく、
同じボタンが押しやすい位置へ並び替わります。ランキング部屋の情報も進捗同期に含まれます。
\end{casebox}

\section*{おすすめの使う順番}

\begin{casebox}[王道ルート：迷ったらこの順番]{forest}
\stepcircle{1}\ \textbf{「証拠を読む」に切り替える}\\
左下の章選択モードを「証拠を読む」にして、最初の事件を開きます。

\stepcircle{2}\ \textbf{学習パートを最後まで読む}\\
物語、完成コード、要点、テストの罠を確認します。途中の確認問題にも答えます。

\stepcircle{3}\ \textbf{「問題パートに切替」}\\
同じ事件の訓練場を開き、最初は低い難易度から挑戦します。

\stepcircle{4}\ \textbf{間違えた問題を「未解決」で回収}\\
間違いは復習ノートへ集まります。正解できるまで短く繰り返します。

\stepcircle{5}\ \textbf{次の事件へ進む}\\
理解してから次へ進むと、後半の複合問題で迷いにくくなります。

\stepcircle{6}\ \textbf{数章終えたら「総合捜査」}\\
1000本ノックで複数単元を混ぜ、本当に思い出せるか確認します。

\stepcircle{7}\ \textbf{ランキング部屋で継続}\\
友達と部屋を作り、普段の進捗を複数の観点から競います。
\end{casebox}

\begin{casebox}[試験直前ルート：時間がないとき]{gold}
\textbf{捜査条件}を開き、\textbf{一夜漬けモード（鉄板・頻出のみ）}を選びます。\\
まず \taglabel{wine}{鉄板}、次に \taglabel{gold}{頻出}。
間違えた問題だけを「未解決」で解き直してください。余裕ができたら標準以上へ進みます。
\end{casebox}

\section*{事件一覧と章攻略}

\begin{casebox}[2つの章選択モード]{navy}
\textbf{証拠を読む：} 解説中心。初見の章、苦手な章、久しぶりの章におすすめ。\\
\textbf{推理に挑む：} 問題中心。理解できているか試したいときにおすすめ。\\
左下のボタンでいつでも切り替えられます。
\end{casebox}

事件カードには章名、内容、評価の星が表示されます。BOSS NODEは複数の知識を組み合わせる
重要事件です。先に通常事件を進めてから挑みましょう。

\subsection*{訓練場の流れ}

\begin{enumerate}
\item 完成コードと解説を確認する。
\item 開始する難易度を選ぶ。
\item 選択式、コード入力、ドラッグ配置などの問題に答える。
\item 正誤と解説を読み、「つづける」で次へ進む。
\item クリア後は復習、再挑戦、次の事件を選ぶ。
\end{enumerate}

\begin{casebox}[すべての問題で「入力／選択」を切替]{navy}
問題上部の切替で、\textbf{入力}または\textbf{選択}を選べます。最初は入力で表示されます。\\
\textbf{入力：} 覚えたコードや用語を自力で思い出す練習。暗記と本番対策に向きます。\\
\textbf{選択：} 候補から答える練習。初見の単元や、ヒントが欲しいときに向きます。\\
問題の正解や進捗はどちらを使っても同じように記録されます。
\end{casebox}

\begin{casebox}[便利なキー操作]{wine}
\key{Enter} 回答・次へ進む \qquad
\key{Esc} 前の画面／地図へ戻る\\
複数行コードでは \key{Enter} で送信、\key{Shift + Enter} で改行します。
\end{casebox}

\section*{難易度と重要度タグ}

\begin{tabularx}{\linewidth}{p{25mm}p{28mm}X}
\textbf{タグ} & \textbf{おすすめ時期} & \textbf{使い方}\\ \hline
\taglabel{wine}{鉄板} & 最初・直前 & 最重要の基礎。ここを落とさない状態を作る。\\
\taglabel{gold}{頻出} & 鉄板の次 & 定番問題。試験前の得点源にする。\\
\taglabel{forest}{標準} & 基礎完成後 & 普通の出題に対応できるか確認する。\\
\taglabel{navy}{深掘り} & 余裕がある時 & 差がつく発展問題に挑戦する。\\
\taglabel{purple}{完全制覇} & 仕上げ & 総復習・高難度。最後の腕試し。\\
\end{tabularx}

\textbf{コツ：} 高難度を無理に選ぶより、鉄板・頻出を高い正答率で終えてから上げるほうが
効率的です。章攻略では、選んだ難易度から上へ自動的にエスカレートします。

\section*{総合捜査（1000本ノック）}

章を横断して問題を解くモードです。「単元」「難易度」「重要度」を組み合わせて出題範囲を
作れます。

\begin{casebox}[おすすめ設定]{forest}
\textbf{初回：} 苦手分野＋おすすめ難易度に自動設定\\
\textbf{試験直前：} 全単元＋鉄板・頻出\\
\textbf{弱点補強：} 苦手な単元だけ＋鉄板から標準\\
\textbf{仕上げ：} 全単元＋深掘り・完全制覇
\end{casebox}

\subsection*{捜査条件で設定できるもの}

\begin{tabularx}{\linewidth}{>{\bfseries}p{32mm}X}
単元 & 全単元、または学習したい単元を複数選択。誤答率も選択の目安になります。\\
重要度 & 鉄板、頻出、標準、深掘り、完全制覇から複数選択できます。\\
難易度 & 初級から上級まで複数選択。全選択ならユーザー未指定として自動調整します。\\
問題数 & \textbf{制限なし（初期値）／10／20／50／100問}から選択します。\\
自動設定 & 苦手分野と、その時点の成績に合うおすすめ難易度を設定します。\\
一夜漬けモード & 全単元の鉄板・頻出だけへ一度に設定します。試験直前向けです。\\
快速モード & 類題をなるべく重ねず、違う内容を優先します。初期値はONです。\\
\end{tabularx}

\begin{casebox}[快速モードの正確な動き]{gold}
快速モードは\textbf{難問モードではありません}。同じ類題タグの2問目以降を、未出の内容より
後ろへ回す機能です。全条件では1,484問が候補で、類題を1問ずつ優先した最初の一巡は
1,422問、後回しになる類題は62問です。問題は削除されず、候補不足や一巡後には再利用します。\\
快速ONだけでは「ユーザーが難易度を指定した」と判定されません。
\end{casebox}

\begin{casebox}[自動難易度と内部レート]{forest}
難易度を指定しない場合、単元ごとの正答・誤答実績から次の問題の難易度を\textbf{1問ごと}に
調整します。問題には表示上の4段階とは別に、より細かな内部レートがあります。
快速モードや問題数を設定しても、この自動調整は止まりません。難易度を明示的に絞った場合だけ、
指定した範囲を優先します。
\end{casebox}

途中で「捜査を中断」しても、ホーム右側から残り問題を再開できます。連続正解、最高記録、
通算正答率を目安にしましょう。

\section*{未解決（復習ノート）}

間違えた問題は自動で集まり、解説画面から手動で復習ノートへ入れることもできます。
「未解決」を開くと、登録された問題だけを集中して解けます。

\begin{casebox}[復習の小ルール]{gold}
一度に全部覚えようとせず、\textbf{間違えた当日 → 翌日 → 数日後}の3回に分けて解き直すと
定着しやすくなります。空になったら、その日の捜査は大成功です。
\end{casebox}

\section*{ランキングルーム}

\begin{enumerate}
\item 「ランキング」を開き、表示名を入力する。
\item 誰か一人が部屋名を決めて「部屋を作成」する。
\item 表示された6文字の招待コードを仲間へ送る。
\item 仲間は「招待コードで参加」から同じ部屋へ入る。
\item 普段どおり学習し、部屋を開いて順位を確認する。
\end{enumerate}

順位は1種類ではなく、総合評価、解決済み事件、学習完了、演習実績など複数の観点で
切り替えられます。進捗はオンラインで自動更新され、「現在の進捗を反映」でも同期できます。

\subsection*{スマホ・PCの進捗同期}

\begin{enumerate}
\item 片方の端末で「ランキング」→「同期コードを発行」を押します。
\item 表示された10文字のコードを、もう片方の端末の「この端末を同期」へ入力します。
\item 学習進捗、設定、参加中のランキング部屋が共有されます。
\item 同期後は短い間隔で自動更新されます。同じコードを知らない端末からは進捗を読めません。
\end{enumerate}

\begin{casebox}[2種類のコードを区別]{navy}
\textbf{6文字の部屋コード：} ランキング部屋へ参加するための招待コード。\\
\textbf{10文字の同期コード：} 自分のスマホとPCで進捗を共有するための秘密のコード。\\
同期コードはパスワードのように扱い、公開場所へ載せないでください。
\end{casebox}

\begin{casebox}[部屋主だけができること]{wine}
部屋を作ったユーザーには「部屋主」と表示されます。部屋主は自分以外のメンバーを
\textbf{キック}できます。「この端末から部屋を削除」は表示から外す操作なので、
使う前に対象の部屋を確認してください。
\end{casebox}

\section*{進捗・実績・共有}

右側の「観測を開く」から、進捗と実績を確認できます。

\begin{itemize}
\item 実績バッジを押すと、取得方法と現在の進み具合を表示。
\item バッジ一覧はページ式。後ろほど高難度の実績。
\item 取得済み実績から5つを進捗状況へドラッグして選択。
\item ユーザー名は進捗状況の中でも変更可能。
\item 進捗状況を画像に保存、テキストをコピー、共有できます。
\item 全取得デモはURL末尾に \texttt{?demo=all} を付けて確認できます。
\end{itemize}

\section*{続けやすくする3つのコツ}

\begin{casebox}[1日5分でも捜査は進む]{forest}
\textbf{1. 今日の一歩だけ終える}\\
大きな目標より、ホームに出る3つの小目標を優先します。

\textbf{2. 連続正解より「戻ってくる」}\\
間違いは未解決リストが覚えてくれます。失敗しても記録が教材になります。

\textbf{3. ランキングは速さではなく継続のきっかけ}\\
得意分野が違っても、ランキングの観点を切り替えれば別の活躍が見えます。
\end{casebox}

\section*{困ったとき}

\begin{tabularx}{\linewidth}{>{\bfseries}p{42mm}X}
表示が古い & ページを再読み込み。改善しない場合はブラウザのキャッシュを更新します。\\
別端末で部屋に入りたい & 同じ招待コードを入力。表示名を確認してから参加します。\\
スマホとPCを同じ進捗にしたい & ランキング画面で10文字の同期コードを発行し、もう一方へ入力します。\\
何から解けばよいか不明 & 「苦手分野＋おすすめ難易度に自動設定」を使います。\\
時間がない & 一夜漬けモードで鉄板・頻出に絞ります。\\
似た問題が多い & 捜査条件の快速モードをONにします。難易度の自動調整は継続します。\\
問題を忘れた & 学習パートまたは訓練場の例題へ戻ります。\\
\end{tabularx}

\vfill
\begin{center}
{\titlejp\Large 捜査は、一問から始まる。}\\[2mm]
{\color{muted}今日の一歩を終えたら、それで十分前進です。}
\end{center}

\end{document}
